\chapter{Literature Review}
\label{chap:literature}

The literature relevant to this thesis was searched between July and
September 2026 using Google Scholar, arXiv, the ACL Anthology, and the
Springer and IEEE digital libraries. The search combined the terms
``automated essay scoring'', ``multi-agent'', ``LLM-as-a-judge'',
``cultural bias'', ``debiasing'', ``fairness'', ``non-native writers'',
and ``contrastive rhetoric''. It was supplemented by backward and forward
citation chasing from key papers. Every reference cited in this thesis
was verified against its primary source. The review is organised
thematically: multi-agent and panel-based scoring
(Section~\ref{sec:lit-multiagent}), inference-time debiasing
(Section~\ref{sec:lit-debias}), cultural bias in LLMs
(Section~\ref{sec:lit-culture}), fairness in AES for non-native writers
(Section~\ref{sec:lit-aes-fairness}), and the contrastive-rhetoric
foundations (Section~\ref{sec:lit-rhetoric}), before
Section~\ref{sec:lit-gap} locates the gap this thesis addresses.

\section{Multi-Agent and Panel-Based Approaches to Scoring}
\label{sec:lit-multiagent}

Two research communities have converged on the same architectural
intuition, that several evaluators are better than one, for
different reasons.

In automated essay scoring, multi-agent designs pursue
\emph{accuracy}. CAFES \citep{su2025cafes} orchestrates scorer,
feedback-manager, and reflective-reviser agents (average 21\% QWK
improvement over direct prompting). Roundtable Essay Scoring
\citep[RES;][]{jang2025roundtable} consolidates trait-based evaluator
agents through simulated discussion (up to 34.86\% on ASAP,
zero-shot). Related designs add debate with retrieval
\citep{keramati2026madrag} and component recognition
\citep{wang2025autoscore}. What unites this family is the axis along
which agents differ: \emph{function} or \emph{trait}. None varies the
cultural interpretation of the rubric itself, and none evaluates
fairness across writer cohorts. Improvements are reported on
aggregate agreement. This leaves open the possibility that a more
accurate panel is uniformly accurate for some writers and uniformly
harsh for others.

In the LLM-evaluation community, panel designs pursue \emph{bias
reduction}. \citet{verga2024replacing} replaced a single large judge
model with a Panel of LLM evaluators (PoLL) drawn from three different
model families. They found that the panel not only correlated better
with human judgment across six datasets but also reduced intra-model
self-preference bias, at roughly a seventh of the cost. This is, to
date, the clearest empirical demonstration that judge \emph{diversity}
counteracts judge \emph{bias}. The diversity axis, however, is model
provider, chosen to decorrelate model-specific quirks rather than to
represent any substantive difference in evaluative standards. The
present thesis can be read as a targeted transplant. It takes the
panel-of-judges mechanism whose bias-reducing effect
\citet{verga2024replacing} established, and re-axes the diversity from
provider to \emph{cultural rubric interpretation}. Here the bias to be
countered is cultural rather than idiosyncratic.

\section{Inference-Time Debiasing}
\label{sec:lit-debias}

A second literature attacks bias at inference time, without retraining.
BiasFilter \citep{cheng2025biasfilter} periodically evaluates candidate
continuations during generation against a fairness reward model and
discards low-reward segments, enforcing fairness on \emph{generated
text} for both open-source and API-based models. The framework
demonstrates that meaningful debiasing is achievable purely at
inference time, a design constraint this thesis shares. But its
unit of intervention is the output token stream, which has no analogue
when the model's output is a single scalar score.

Closest to the present work is Multi-Agent Cultural Debate (MACD)
\citep{tan2026macd}, which assigns LLM agents explicit cultural
personas and orchestrates deliberation. MACD substantially raised the
rate of unbiased responses over single-model baselines, and its
central finding, that functionally-roled agent frameworks lack the
explicit cultural representation cross-cultural fairness needs,
independently corroborates the design principle behind CAMPS. The two
frameworks nonetheless answer different questions. MACD debiases
\emph{what a model says} in open-ended generation. CAMPS debiases
\emph{how a model judges}: the deliverable is a score whose
fairness can be measured cohort by cohort, and disagreement
between cultural agents is exploited as a detection signal rather than
dissolved through debate.

A caution for all multi-agent designs comes from the MALIBU benchmark
\citep{mirza2025malibu}, which shows that LLM-based multi-agent systems
can implicitly \emph{reinforce} social biases, and that
persona-conditioned judging can over-correct. For this thesis, MALIBU
motivates two design decisions: consensus weights are calibrated
empirically rather than assumed benign (Hypothesis H4 measures the
over-correction frontier directly), and panel behaviour is validated
against multi-rater human ground truth rather than taken on faith.
A separate line of work mitigates bias inside the model rather than
around it. Self-debiasing exploits a model's own ability to recognise
undesirable properties of its output, suppressing them at decoding
time \citep{schick2021self}, and FairSteer intervenes on hidden-layer
activations with dynamically applied steering vectors at inference
time \citep{li2025fairsteer}. Both are effective in their settings,
but both are white-box methods: they require access to output
distributions or internal activations that commercial scoring
deployments accessed through an API do not expose. CAMPS was
deliberately designed to operate at the level of prompts and outputs
alone, trading the precision of internal intervention for
applicability to any evaluator, including closed models.

\section{Cultural Bias in Large Language Models}
\label{sec:lit-culture}

That LLMs carry a Western-centric evaluative disposition is no longer
seriously contested. \citet{tao2024cultural} compared five OpenAI model
generations against nationally representative values-survey data for
107 countries. They found that model responses consistently resemble
those of English-speaking and Protestant-European populations. This is
the WEIRD alignment that this thesis takes as the distal cause of
evaluative epistemic filtering.

The proximal framework for this thesis comes from the supervisor's
research programme. \citet{dorleon2026uncovering} formally defined
bias categories for generative outputs (stereotyping, omission,
ethnocentrism, simplification) and proposed multi-perspective
generation with bias detection, proving that under ideal conditions the
method eliminates stereotypical content and perspective omission. The
companion study \citep{shujaa2026llms} evaluated seven state-of-the-art
LLMs on prompts embedding culturally sensitive assumptions about
Vietnam, Myanmar, and Nepal, documenting pervasive stereotyping,
linguistic ethnocentrism, and epistemic bias, with substantial
variation across models. It is this cross-model variation that
motivates testing H1 on at least four model families. The same group's work on
framing-bias detection \citep{shujaa2025framing} reflects the
detection-oriented strand of this programme that the Bias Divergence
Score continues: bias treated not only as something to remove but as
something to \emph{measure and flag}.
Two further studies shaped how cultural calibration is implemented in
this thesis. SEA-HELM, a community-built evaluation suite for
Southeast Asian languages, argues that benchmarks produced by machine
translation or synthetic generation without community input lose
cultural authenticity \citep{susanto2025seahelm}. The same reasoning
led this thesis to ground its cultural lenses in the contrastive
rhetoric literature and in a corpus built in the region, rather than
in ad hoc persona descriptions. More cautionary still,
\citet{khan2025randomness} showed that survey-based measures of LLM
cultural alignment are unstable across presentation formats and that
prompting a model to represent a specific culture is not reliably
steerable. For CAMPS this is a design constraint, not a refutation:
lens prompts are fixed verbatim (Appendix~A), decoding is
deterministic, and the panel's behaviour is validated empirically
against human raters (E3) rather than assumed correct because the
prompt requested it.

\section{Fairness in AES for Non-Native Writers}
\label{sec:lit-aes-fairness}

Concerns about equitable automated scoring of non-native writing
predate the LLM era. \citet{vajjala2018automated} showed on the
TOEFL11 and FCE corpora that the linguistic features most predictive of
essay quality differ across L1 populations, evidence that a single
scoring model implicitly encodes population-specific assumptions about
what good writing looks like. What has changed with LLM-based scoring
is the depth at which such assumptions operate: where feature-based
systems encoded them in engineered predictors that could be inspected,
LLM evaluators encode them in an opaque prior over rhetorical form
\citep[cf.][]{gallegos2024bias}.

The LLM-era evidence is accumulating quickly.
\citet{pack2024llm} evaluated four prominent LLMs (PaLM~2, Claude~2,
GPT-3.5, GPT-4) as scorers of English-learner placement essays,
finding usable reliability for the strongest model but meaningful
variation across models and administrations. Evaluative behaviour is
therefore model-specific, which motivates testing H1
across at least four model families rather than generalising from one.
Most directly relevant, \citet{gayed2026firstlanguage} fine-tuned an
open-weight model (Gemma-3-27B) for essay scoring and evaluated it on
the full TOEFL11 corpus of 12{,}100 essays across eleven L1
backgrounds, explicitly analysing scoring effects by first language.
The study documents L1-linked scoring variation on exactly the corpus
family and task that concern this thesis. Like the broader AES
fairness literature it extends, however, it stops at \emph{measurement}: no
mitigation framework is proposed, and the native/non-native design
this thesis requires is out of scope because TOEFL11 contains no
native-English cohort. Related recent work has begun to probe
mechanism as well as outcome. \citet{yang2025demographics} showed that
prompt-based LLM scorers can infer writer demographics, including
English-learner status, from essay text alone. They also showed that
scoring behaviour shifts once such a status is recognised.

Taken together, this literature establishes the disparity, its
model-dependence, and hints at its mechanism, while leaving the
mitigation side of the problem essentially untouched. That asymmetry
between diagnosis and treatment is the practical motivation for the
framework developed in Chapter~3.

\section{Contrastive Rhetoric and Cultural Writing Traditions}
\label{sec:lit-rhetoric}

The claim that rhetorical organisation is culturally learned, and
that an evaluator calibrated to one tradition will therefore
systematically misread another, has a research lineage of its own.
\citet{kaplan1966cultural} first argued that paragraph organisation
reflects culturally specific ``thought patterns'', contrasting the
linear development expected in English academic prose with the
patterns he observed in other traditions. The framework was refined
substantially in later work: \citet{hinds1987reader} distinguished
\emph{writer-responsible} languages, in which the burden of clarity
falls on the writer through explicit signposting, from
\emph{reader-responsible} languages, in which coherence is expected to
be inferred by the reader. Under this typology, the ``missing
thesis statement'' an LLM penalises may simply be reader-responsible
coherence operating as designed. \citet{connor2011intercultural}
consolidated this field as intercultural rhetoric, correcting early
essentialist readings: rhetorical traditions are conventions
distributed across writing communities, not fixed properties of
individuals. This thesis adopts that distinction explicitly, treating
L1 cohorts as proxies for rhetorical tradition rather than for
identity.

For Southeast Asian learner writing specifically, the empirical base
is anchored by the ICNALE project \citep{ishikawa2023icnale}, whose
topic-controlled design allows rhetorical variation to be observed
under matched task conditions across ten Asian countries and a
native-English reference cohort. This literature performs two duties in
the present thesis. First, it supplies the content of the cultural rubric
interpretations in Chapter~3, where each CRI prompt is grounded in
documented rhetorical conventions rather than improvised cultural
personas. Second, it supplies the defence against the objection that
culturally-calibrated agents merely role-play stereotypes. The
calibrations are drawn from a half-century of contrastive-rhetoric
scholarship, and their behaviour is tested empirically in Chapter~4.

\section{The Research Gap}
\label{sec:lit-gap}

Three literatures converge on the problem this thesis addresses, and
each stops short of it in a characteristic way. The multi-agent scoring
literature (Section~\ref{sec:lit-multiagent}) established that panels
of evaluators outperform single judges (CAFES and RES on accuracy,
PoLL on bias). But it defines panel diversity by function, trait, or
model provider, never by cultural interpretation of the rubric, and it
does not measure fairness across writer cohorts. The inference-time
debiasing literature (Section~\ref{sec:lit-debias}) established that
bias can be mitigated without retraining, and MACD specifically
established that \emph{explicit cultural representation} in agent
frameworks matters. But its object is generated text, not evaluative
judgment, and it dissolves inter-agent disagreement rather than
exploiting it. The cultural-bias literature
(Section~\ref{sec:lit-culture}) established the WEIRD disposition of
current models and, in \citet{dorleon2026uncovering}, provided a formal
multi-perspective mitigation framework with provable guarantees, for
generation only. Meanwhile, the AES fairness literature
(Section~\ref{sec:lit-aes-fairness}) continues to document L1-linked
scoring disparities without offering a validated mitigation framework.

The intersection itself remains empty: culturally-diverse rubric
interpretation, applied to evaluative scoring, with inter-agent
divergence used as a calibrated trigger for human review and validated
against multi-rater human ground truth. To the best of the searches
described above, no published work occupies it.
That intersection is precisely where CAMPS operates, and the recency of
the nearest neighbours on either side suggests the space will not
remain unoccupied for long.

\section*{Chapter summary}
Two literatures converge on the idea that several evaluators are
better than one, but neither varies the \emph{cultural} reading of
the rubric, and the fairness literature documents L1-linked scoring
disparities without a validated remedy. Chapter~3 builds the
framework that fills this gap and the experiments that test it.
